% Libro 3 · Mundo 4 · Vocabulario en contexto (app: pu-vocabulario)
\mundo{Vocabulario en contexto}{pu-vocabulario}{En el examen, la pregunta no es «¿qué significa esta palabra?» sino «¿qué
significa \emph{en este texto}?». Una misma palabra cambia de sentido según la frase que la rodea, y el diccionario de
memoria no basta: hay que leer el contexto. Este mundo enseña a hacerlo con sinónimos, palabras de varios sentidos,
expresiones y las piezas con que se forman las palabras.}{%
\item \textbf{4.1} Sinónimos según el contexto
\item \textbf{4.2} Palabras con varios sentidos
\item \textbf{4.3} Expresiones y frases hechas
\item \textbf{4.4} Prefijos y raíces}

% ------------------------------------------------------------------
\leccion{pu-vo-1}{Sinónimos según el contexto}{Juego de cambiar una palabra por otra sin romper la frase.}

\begin{reto}
«La propuesta fue \emph{acogida} por la mayoría de los vecinos.» En esta oración, \emph{acogida} puede reemplazarse por
\opcionesd{hospedada}{aceptada}{protegida}{recogida}
\end{reto}

\tecnica{La técnica: la prueba del reemplazo}
Para elegir un sinónimo en contexto, se hace la \clave{prueba del reemplazo}: se pone cada opción en lugar de la palabra y
se lee la oración completa. La correcta conserva el sentido \emph{y} suena natural. Los distractores suelen ser sinónimos
reales de la palabra, pero en otro de sus sentidos: «acoger» puede ser «hospedar» (acoger a un viajero), pero una
propuesta no se hospeda.

\begin{regla}
\begin{enumerate}
  \item Antes de mirar las opciones, di con tus palabras qué significa la palabra en esa oración.
  \item Reemplaza y lee la oración entera con cada opción.
  \item Descarta los sinónimos de otro sentido de la palabra.
\end{enumerate}
\end{regla}

\begin{texto}{Elogio de la lentitud}{Ensayo · texto original escrito para este libro}
\pa Vivimos rodeados de promesas de rapidez: comida en diez minutos, idiomas en treinta días, títulos en un año. La
velocidad se ha vuelto una virtud tan obvia que casi nadie la \emph{cuestiona}. Y, sin embargo, algunas de las cosas que
más valoramos solo crecen despacio.

\pa Un árbol no se apura porque lo miremos con impaciencia. Una amistad tampoco: se \emph{fragua} en años de
conversaciones, favores y silencios compartidos. Hasta el pan, que hoy se vende en minutos, necesita horas de reposo para
que la masa \emph{leve}.

\pa Lo mismo ocurre con el aprendizaje. Quien estudia a última hora puede aprobar un examen, pero lo aprendido así es
\emph{efímero}: se evapora a los pocos días. Lo que se aprende con calma, en cambio, \emph{arraiga}, y sirve para aprender
lo siguiente.

\pa No se trata de hacer todo despacio. Hay urgencias reales, y hay lentitudes que son solo pereza. Se trata de
\emph{discernir}: saber qué cosas ganan con la prisa y cuáles se arruinan con ella.
\end{texto}

\subsection{Preguntas comentadas}
\begin{comentada}{En el \S\,3, \emph{efímero} significa}
\mal{A}{difícil}{No tiene relación con la dificultad.}
\bien{B}{pasajero, que dura poco}{El propio texto lo explica: «se evapora a los pocos días».}
\mal{C}{falso}{\emph{Dice otra cosa}: lo aprendido no es falso, sino breve.}
\mal{D}{importante}{Contradice el sentido de la frase.}
\end{comentada}

\taller
\piso{Como en el examen}{opción múltiple}
\begin{ejercicios}
\item En el \S\,1, \emph{cuestiona} puede reemplazarse por
\opcionesd{pregunta}{pone en duda}{responde}{estudia}
\sol{B. «Casi nadie la pone en duda».}
\item En el \S\,2, la amistad «se \emph{fragua}» significa que
\opcionesl{se destruye.}{se va formando poco a poco.}{se calienta.}{se olvida.}
\sol{B. Como el metal en la fragua, se forma con tiempo.}
\item En el \S\,3, \emph{arraiga} significa
\opcionesd{echa raíces, se afianza}{se arranca}{se repite}{se olvida}
\sol{A. Contrasta con «efímero».}
\item En el \S\,4, \emph{discernir} significa
\opcionesd{distinguir con criterio}{apresurarse}{descansar}{discutir}
\sol{A.}
\end{ejercicios}

\begin{retoresuelto}
Una propuesta «acogida» por los vecinos es una propuesta que ellos aceptaron. «Hospedada» y «protegida» son sentidos de
\emph{acoger} que no encajan aquí. \textbf{Opción B.}
\end{retoresuelto}
\bitacora

% ------------------------------------------------------------------
\leccion{pu-vo-2}{Palabras con varios sentidos}{Juego de la palabra que cambia de traje.}

\begin{reto}
«El gerente presentó un \emph{balance} positivo del año.» Aquí, \emph{balance} significa
\opcionesd{equilibrio físico}{movimiento de un lado a otro}{evaluación de resultados}{columpio}
\end{reto}

\tecnica{La técnica: el contexto elige el sentido}
Muchas palabras son \clave{polisémicas}: tienen varios sentidos relacionados (\emph{balance}: equilibrio, cuentas,
evaluación). Otras son \clave{homónimas}: suenan o se escriben igual pero vienen de palabras distintas (\emph{vino} del
verbo venir y \emph{vino} la bebida). En ambos casos, lo que decide es el contexto: el tema del texto, las palabras
vecinas, lo que se dice antes y después.

\subsubsection{El sentido figurado}
Muchas palabras tienen un sentido \emph{literal} (físico, concreto) y otro \emph{figurado} (abstracto): una «puerta» de
madera y «la puerta» a una oportunidad; una «raíz» de árbol y «la raíz» de un problema.

\begin{regla}
\begin{itemize}
  \item Identifica el tema del párrafo: decide en qué «terreno» está la palabra.
  \item Mira las palabras vecinas: «balance» junto a «del año» y «positivo» es una evaluación.
  \item Distingue el sentido literal del figurado.
\end{itemize}
\end{regla}

\begin{texto}{Las palabras de la finca}{Crónica breve · texto original escrito para este libro}
\pa Mi abuelo tenía una finca en la \emph{sierra} y una \emph{sierra} en el galpón. Con la primera se ganaba la vida; con la
segunda, la perdía un poco cada vez que la usaba sin guantes. De niño me costaba entender que una misma palabra sirviera
para una montaña y para una herramienta, hasta que él me explicó que las dos tenían dientes.

\pa Los domingos bajaba al pueblo a vender café. Llegaba al \emph{banco} de la plaza, se sentaba a esperar que abrieran el
otro \emph{banco}, el de la plata, y mientras tanto conversaba con todo el que pasaba. Decía que en el primero se ganaba más
que en el segundo, porque allá se cobraban intereses y aquí se regalaban noticias.

\pa Cuando murió, el pueblo entero fue al entierro. El cura dijo que había dejado una \emph{huella} profunda, y yo, que
tenía nueve años, pensé en las botas de caucho que dejaban marcas en el barro del camino. Tardé años en entender que el
cura y yo hablábamos de lo mismo.
\end{texto}

\subsection{Preguntas comentadas}
\begin{comentada}{En el \S\,3, la «huella profunda» que menciona el cura se refiere a}
\mal{A}{las marcas de las botas en el barro.}{Es lo que pensó el niño: el sentido \emph{literal}.}
\bien{B}{el recuerdo e influencia que el abuelo dejó en la gente.}{Sentido figurado; el narrador dice que tardó en
entenderlo.}
\mal{C}{una herida del abuelo.}{\emph{No lo dice}.}
\mal{D}{el camino hacia el pueblo.}{\emph{Dice otra cosa}.}
\end{comentada}

\taller
\piso{Como en el examen}{opción múltiple}
\begin{ejercicios}
\item Según el abuelo, la sierra-montaña y la sierra-herramienta se parecen en que
\opcionesl{las dos son peligrosas.}{las dos «tienen dientes».}{las dos están en la finca.}{las dos se usan sin guantes.}
\sol{B. \S\,1: los picos de la montaña y los dientes de la herramienta.}
\item En el \S\,2, el abuelo opina que en el banco de la plaza «se ganaba más» porque allí
\opcionesl{le prestaban dinero.}{recibía noticias y conversación gratis.}{vendía más café.}{no cobraban intereses por los créditos.}
\sol{B.}
\item La última oración del texto indica que el narrador
\opcionesl{nunca entendió al cura.}{con el tiempo comprendió que la huella literal y la figurada estaban relacionadas.}{olvidó a su abuelo.}{no estuvo en el entierro.}
\sol{B.}
\end{ejercicios}

\begin{retoresuelto}
En el contexto de un gerente y de los resultados del año, \emph{balance} es la evaluación de lo logrado. \textbf{Opción C.}
\end{retoresuelto}
\bitacora

% ------------------------------------------------------------------
\leccion{pu-vo-3}{Expresiones y frases hechas}{Juego de unir la expresión con lo que de verdad quiere decir.}

\begin{reto}
«Después de tres intentos fallidos, Mateo \emph{tiró la toalla}.» Esto significa que Mateo
\opcionesd{se secó}{se rindió}{se enojó}{se fue a nadar}
\end{reto}

\tecnica{La técnica: no leer palabra por palabra}
Una \clave{frase hecha} (o locución) tiene un significado que no sale de sumar sus palabras: «tirar la toalla» no tiene
nada que ver con toallas. Muchas vienen de oficios o costumbres (el boxeo, el campo, el mar). En el examen, si una opción
interpreta la expresión literalmente, casi siempre es un distractor.

\begin{center}\small
\begin{tabularx}{\linewidth}{@{}l>{\raggedright\arraybackslash}X@{}}
\toprule
\textbf{Expresión} & \textbf{Sentido}\\\midrule
meter la pata & equivocarse, decir algo inoportuno\\
dar en el clavo & acertar\\
estar entre la espada y la pared & tener que elegir entre dos opciones malas\\
llover sobre mojado & sumarse un problema a otro\\
echar leña al fuego & empeorar un conflicto\\
dar gato por liebre & engañar dando algo de menor valor\\
a regañadientes & de mala gana\\\bottomrule
\end{tabularx}
\end{center}

\begin{regla}
\begin{itemize}
  \item Si una frase no tiene sentido literal en el contexto, es una expresión.
  \item Deduce el sentido por la situación: ¿qué pasó antes y después?
  \item Descarta las opciones literales.
\end{itemize}
\end{regla}

\begin{texto}{La asamblea del edificio}{Relato · texto original escrito para este libro}
\pa La reunión empezó tarde y mal. El administrador anunció que había que pintar la fachada y que la cuota subiría, y en el
salón comunal \emph{se armó la de Troya}. El señor del 302 dijo que eso era \emph{llover sobre mojado}, porque el año anterior
ya habían pagado el arreglo del ascensor.

\pa Doña Ofelia, que nunca habla, pidió la palabra y \emph{dio en el clavo}: propuso pintar solo la parte de la fachada que
da a la calle y dejar el resto para el año siguiente. Algunos aplaudieron; el administrador aceptó \emph{a regañadientes}.

\pa Todo iba bien hasta que el del 501 \emph{echó leña al fuego} recordando que el administrador cobraba un sueldo «por no
hacer nada». La discusión se encendió otra vez y la reunión terminó a medianoche, sin votación.
\end{texto}

\subsection{Preguntas comentadas}
\begin{comentada}{En el \S\,1, «se armó la de Troya» significa que}
\mal{A}{se organizó una fiesta.}{Lectura literal y equivocada.}
\bien{B}{se produjo un gran alboroto.}{Coherente con la noticia de la cuota y las quejas que siguen.}
\mal{C}{se habló de historia.}{Literal.}
\mal{D}{todos estuvieron de acuerdo.}{\emph{Dice lo contrario}.}
\end{comentada}

\taller
\piso{Como en el examen}{opción múltiple}
\begin{ejercicios}
\item Que doña Ofelia «dio en el clavo» quiere decir que
\opcionesd{golpeó la mesa}{hizo una propuesta acertada}{se equivocó}{se quejó}
\sol{B.}
\item El del 302 dice que es «llover sobre mojado» porque
\opcionesl{estaba lloviendo.}{el nuevo cobro se sumaba a otro reciente.}{el ascensor estaba mojado.}{la pintura se iba a mojar.}
\sol{B.}
\item Cuando el del 501 «echó leña al fuego»,
\opcionesl{calmó la discusión.}{empeoró la discusión.}{propuso una solución.}{se fue de la reunión.}
\sol{B. «La discusión se encendió otra vez».}
\end{ejercicios}

\begin{retoresuelto}
En el boxeo, el entrenador tira la toalla para detener la pelea. Tras tres intentos fallidos, Mateo se rindió.
\textbf{Opción B.}
\end{retoresuelto}
\bitacora

% ------------------------------------------------------------------
\leccion{pu-vo-4}{Prefijos y raíces}{Juego de armar palabras con piezas.}

\begin{reto}
Si \emph{omni-} significa «todo» y \emph{-voro} «que come», un animal \emph{omnívoro} es el que
\opcionesl{come solo plantas.}{come de todo.}{no come nada.}{come una sola vez al día.}
\end{reto}

\tecnica{La técnica: desarmar la palabra}
Muchas palabras del español están hechas de piezas: un \clave{prefijo} (al principio), una \clave{raíz} (el núcleo) y un
\clave{sufijo} (al final). Conocer las piezas más frecuentes permite adivinar el significado de una palabra nunca vista.

\begin{center}\small
\begin{tabular}{@{}ll@{\qquad}ll@{}}
\toprule
\textbf{Pieza} & \textbf{Sentido} & \textbf{Pieza} & \textbf{Sentido}\\\midrule
a-, an- & sin, no (\emph{anónimo}) & -logía & estudio (\emph{biología})\\
anti- & contra (\emph{antídoto}) & -fobia & miedo, rechazo\\
bi- & dos (\emph{bilingüe}) & -cida & que mata (\emph{insecticida})\\
hiper- & exceso (\emph{hipertensión}) & -cracia & gobierno (\emph{democracia})\\
hipo- & debajo, poco (\emph{hipotermia}) & bio- & vida\\
pre- & antes (\emph{prever}) & geo- & tierra\\
re- & de nuevo (\emph{reconstruir}) & grafo, grafía & escritura\\
sub- & debajo (\emph{subterráneo}) & tele- & lejos\\\bottomrule
\end{tabular}
\end{center}

\begin{regla}
\begin{itemize}
  \item Separa prefijo, raíz y sufijo; traduce cada pieza.
  \item Comprueba el sentido armado con el contexto de la oración.
  \item Cuidado con los falsos amigos: \emph{anillo} no lleva el prefijo «a-».
\end{itemize}
\end{regla}

\begin{texto}{Palabras que nacen}{Divulgación · texto original escrito para este libro}
\pa Cada año aparecen en el español palabras nuevas. Algunas llegan de otros idiomas, como \emph{wifi}; otras se forman con
piezas que ya teníamos. Cuando alguien dijo por primera vez \emph{teletrabajo}, nadie tuvo que explicar mucho: \emph{tele-}
ya significaba «a distancia» en \emph{teléfono} y \emph{televisión}.

\pa Este mecanismo es antiguo. Hace siglos, los científicos tomaron raíces griegas y latinas para nombrar lo que
descubrían: \emph{microscopio} (ver lo pequeño), \emph{hidrografía} (escritura o descripción de las aguas),
\emph{biodiversidad} (variedad de la vida). Por eso un estudiante que conoce veinte raíces puede descifrar cientos de
términos técnicos.

\pa Las palabras nuevas no siempre son bien recibidas. Algunos hablantes las consideran innecesarias o feas; otros
celebran que la lengua se mueva. Lo cierto es que el idioma, como cualquier ser vivo, solo deja de cambiar cuando muere.
\end{texto}

\subsection{Preguntas comentadas}
\begin{comentada}{Según el texto, \emph{hidrografía} significa}
\mal{A}{miedo al agua}{Eso sería \emph{hidrofobia}.}
\bien{B}{descripción de las aguas}{El texto la explica entre paréntesis: «escritura o descripción de las aguas».}
\mal{C}{agua pura}{\emph{No lo dice}.}
\mal{D}{estudio de la tierra}{Sería \emph{geología}.}
\end{comentada}

\taller
\piso{Calienta}{desarmar}
\begin{ejercicios}
\item ¿Qué significan \emph{subacuático}, \emph{prehistoria} y \emph{bisemanal}? \sol{Debajo del agua; antes de la historia (escrita); dos veces por semana.}
\end{ejercicios}
\piso{Como en el examen}{opción múltiple}
\begin{ejercicios}
\item Según el \S\,1, \emph{teletrabajo} se entendió fácilmente porque
\opcionesl{venía del inglés.}{el prefijo \emph{tele-} ya era conocido en otras palabras.}{era una palabra muy antigua.}{lo explicó la Academia.}
\sol{B.}
\item La afirmación final, «el idioma solo deja de cambiar cuando muere», expresa que
\opcionesl{el cambio es una señal de vida de la lengua.}{las palabras nuevas matan el idioma.}{el español está muriendo.}{hay que prohibir las palabras nuevas.}
\sol{A.}
\item Una persona \emph{políglota} es la que
\opcionesd{habla muchas lenguas}{estudia a los políticos}{tiene muchos dientes}{escribe poesía}
\sol{A. \emph{Poli-}, muchos; \emph{-glota}, lengua.}
\end{ejercicios}

\begin{retoresuelto}
\emph{Omni-} (todo) $+$ \emph{-voro} (que come): que come de todo. \textbf{Opción B.}
\end{retoresuelto}
\bitacora

% ------------------------------------------------------------------
\evaluacion
\begin{texto}{El oficio de remendar}{Ensayo breve · texto original escrito para este libro}
\pa Remendar fue durante siglos una tarea \emph{cotidiana}. Una camisa rota no se botaba: se le ponía un parche, se zurcía
el codo, se cambiaba el cuello gastado. Lo hacían las abuelas al anochecer, con la paciencia de quien sabe que nada se
consigue fácil.

\pa Hoy la ropa es tan barata que repararla parece una \emph{excentricidad}. Comprar una prenda nueva cuesta menos que el
tiempo de arreglar la vieja, y así los armarios se llenan y los basureros también. La industria textil es, según diversos
estudios, una de las que más agua consume y más desechos produce.

\pa Por eso, en algunas ciudades, el remiendo ha vuelto, pero con otro \emph{talante}: ya no por necesidad sino por
convicción. Hay talleres donde los jóvenes aprenden a coser parches de colores que no esconden el daño, sino que lo
exhiben. La cicatriz, en lugar de ser una vergüenza, se convierte en adorno.

\pa Quizás ahí esté la lección. Remendar no es solo ahorrar: es aceptar que las cosas, como las personas, pueden romperse
y seguir sirviendo. En una época que \emph{desecha} con tanta facilidad, aprender a reparar es casi un acto de rebeldía.
\end{texto}
\begin{ejercicios}
\item En el \S\,1, \emph{cotidiana} significa
\opcionesd{diaria, habitual}{difícil}{antigua}{costosa}
\sol{A. (4.1)}
\item En el \S\,2, \emph{excentricidad} se refiere a algo
\opcionesd{raro, fuera de lo común}{necesario}{ilegal}{central}
\sol{A. (4.1)}
\item En el \S\,3, \emph{talante} puede reemplazarse por
\opcionesd{talento}{actitud, ánimo}{tela}{precio}
\sol{B. (4.2)}
\item La frase «La cicatriz se convierte en adorno» \pf{3} se refiere, en sentido figurado, a
\opcionesl{las heridas de las costureras.}{los parches visibles que muestran el arreglo con orgullo.}{la moda de hacerse tatuajes.}{la ropa nueva.}
\sol{B. (4.2)}
\item «Un acto de rebeldía» \pf{4}, en el contexto, sugiere que reparar
\opcionesl{es ilegal.}{va en contra de la costumbre dominante de desechar.}{es una protesta violenta.}{lo hacen solo los jóvenes.}
\sol{B. (4.3)}
\item En \emph{desecha} \pf{4}, el prefijo \emph{des-} indica
\opcionesd{negación o separación}{repetición}{exceso}{anterioridad}
\sol{A. (4.4)}
\end{ejercicios}
